\chapter{数据，问题的载体：数据思维入门}

数据不是现实本身，而是现实经过测量、筛选、编码和保存之后留下的记录。这个区别构成数据思维的起点：模型只能学习记录中存在的规律，而无法自动恢复从未被采集、被制度性忽略或在处理过程中丢失的信息。高风险 AI 项目的研究发现，许多后续模型故障并非源于算法，而是源于上游数据定义和组织问题沿流程逐级放大形成的“数据级联”\cite{sambasivan2021data}。

数据是机器学习项目的物质基础，但它从来不是现实世界的透明镜像。每一份数据集都带着采集方式的烙印：检测器的物理特性、调查问卷的措辞方式、平台的派单逻辑——这些都深刻地影响着数据记录的内容和形态。不理解数据的来历，就无法判断数据的边界；不判断数据的边界，就无法设计出可靠的模型。

本章讨论数据思维的核心问题：数据是什么、一行数据代表什么、标签是怎么产生的、数据质量问题如何影响分析、以及如何向 AI 描述你的数据。本章不涉及数据清洗的具体操作（第五章）和特征构造（第六章），而是聚焦于在动手处理数据之前，你需要建立的基本认知框架。

% ----------------------------------------------------------------
\section{数据是什么：从现象到记录的映射}
% ----------------------------------------------------------------

数据不是现实的直接复制，而是现实经过特定采集过程之后留下的痕迹。这个区别看起来显而易见，但在实际工作中经常被忽视，并由此引发一系列建模错误。

\subsection{采集过程决定数据含义}

同名字段只有在采集机制相同的情况下才具有可比性。例如“速度”可能来自线圈检测器的点速度、浮动车的区间速度，也可能来自人工填报；它们的误差结构、覆盖对象和缺失规律完全不同。理解数据因此不能止于读取字段类型，而要追问谁在何时、出于什么目的、使用什么设备或规则记录了它。为数据集建立标准化说明文档的“Datasheets for Datasets”工作，正是试图把动机、构成、采集、预处理、用途和维护责任显式化\cite{gebru2021datasheets}。

同一个交通现象，用不同的方式记录，得到的数据在含义上可以相差甚远。以“某路段的车辆流量”为例：感应线圈检测器计的是压过线圈的轴次，再折算成辆数，但在多车道场景中存在车道间干扰；视频检测器依靠图像识别算法计数，在遮挡、光线不足或拥挤场景下识别率下降；浮动车（出租车、网约车）GPS 数据只能捕捉携带 GPS 设备的车辆，而这部分车辆在全体车辆中的占比因时因地而异；手机信令数据依赖基站覆盖密度和用户开机状态，在地下隧道或信号盲区中会出现系统性缺失。

因此，“车流量”这个看似统一的概念，在不同数据集中实际上对应着不同的操作定义。把来自不同采集方式的数据混用，而不清楚各自的含义，是一种非常常见的分析错误。

\subsection{交通领域主要数据类型及其内生局限}

交通工程研究和实践中常见的数据类型，各自有独特的采集机制和由此产生的内生局限。

\textbf{固定检测器数据}包括感应线圈、地磁传感器、视频和毫米波雷达等固定设备，安装在特定位置，以固定频率汇总数据（通常为1分钟或5分钟一次），优点是覆盖时间长、成本摊薄后较低。主要局限在于只能覆盖安装位置附近，全路网覆盖率通常不高；设备会老化漂移，需要定期标定；检测器故障时，数据可能变为零、负数或固定值，而不是缺失，容易被当作正常数据处理。

\textbf{GPS 轨迹数据}来自车载 GPS、手机导航 APP 和共享单车/汽车平台，提供高时空分辨率的移动信息。主要局限在于城市峡谷效应导致定位漂移；隧道、地下停车场等遮挡环境中信号中断，轨迹出现“跳变”；轨迹只代表持有设备的出行群体，存在系统性的人群偏差（例如网约车轨迹不代表私家车行为）。

\textbf{公交刷卡与 APP 数据}记录公交、地铁的 IC 卡刷卡和乘车 APP 数据，包含上下车时间、站点和支付信息。主要局限在于仅能观测到持卡或使用 APP 的乘客，现金支付或逃票行为无法被记录；IC 卡系统通常只记录上车刷卡，下车信息需要通过换乘记录或末站推断；APP 数据受到平台覆盖率的限制，不同城市差异较大。

\textbf{问卷调查数据}来自居民出行调查（如住户调查、网络问卷），提供出行目的、方式选择和社会经济属性等固定检测器无法获取的信息。主要局限在于样本量受成本限制，通常无法覆盖大规模人群；回忆偏差（受访者对过去出行的记忆不准确）和社会期望偏差（受访者倾向于给出“正确答案”）普遍存在；调查频率低，无法反映动态变化。

\textbf{事故记录数据}来自警方或保险公司，包含事故时间、地点、严重程度、涉及车辆等信息。主要局限在于严重依赖报告率，轻微事故漏报率极高；记录质量受警察局、派出所的执行标准影响，同一地区不同时间段的记录详细程度可能不同；地点信息通常为文字描述，映射到地图坐标存在误差。

\textbf{遥感与航拍数据}通过卫星图像和无人机航拍覆盖大范围区域，捕捉停车占用、路面状况和人流聚集等地面传感器难以监测的现象。主要局限在于受云层、光照条件影响，部分时间段无法成像；时间分辨率较低（商业卫星重访周期通常为数天）；图像解析需要算法识别，识别精度受训练数据质量影响。

\textbf{地图与 API 数据}来自高德、百度、Google Maps 等平台，提供路网信息、实时路况、兴趣点（POI）和出行时间估算。主要局限在于路况数据来自平台自身的浮动车和用户上报，覆盖率因地区和时段而异；API 返回的“预计行程时间”是平台算法的估算结果，不是直接测量值；POI 数据的更新频率和完整性因地区不同差异显著。

\subsection{核心洞察：数据是带有偏见的记录}

上述例子揭示了一个核心结论：\textbf{每一份数据都有其采集逻辑，而采集逻辑决定了哪些现象被记录、被谁记录、以何种精度记录}。在分析之前，理解采集逻辑是理解数据含义的前提。

这个认识对 AI 协作具有直接含义：AI 看到的是字段名和数值，看不到这些数值背后的采集机制。如果你不向 AI 说明采集逻辑，它会用对数据含义的通用假设来填补这个空白——而这些假设在你的具体场景中很可能是错误的。

% ----------------------------------------------------------------
\section{一行数据代表什么：观测单位的重要性}
% ----------------------------------------------------------------

理解数据的第一步，不是看字段名，而是搞清楚：\textbf{数据集中的每一行代表什么}。这个问题看似简单，但在实际工作中，忽视观测单位是引发建模错误的高频原因之一。

\subsection{观测单位的概念}

观测单位不仅是表格设计问题，还决定了统计推断的对象。整洁数据原则提出，每个变量构成一列、每个观测构成一行、每类观测单位构成一张表\cite{wickham2014tidy}。这一原则有助于组织分析，但“一行应该代表什么”仍须由研究问题决定：同一批刷卡记录既可整理为一次乘车、一名乘客的一天，也可聚合为一个站点的一小时；每种选择都会保留某些关系并丢失另一些关系。

观测单位（unit of observation）是指数据集中每一条记录所对应的实体或事件。不同的观测单位意味着完全不同的建模设置。交通领域常见的观测单位类型各有其适用场景。\textbf{车辆-时刻}单位中每行是某辆车在某个时刻的状态（如车速、位置），来自 GPS 轨迹或车载传感器，适合分析单车行为，但同一辆车在不同时刻的记录存在强时序依赖。\textbf{路段-时间片}单位中每行是某路段在某个时间段（如5分钟、1小时）内的汇总统计（如平均速度、总流量），来自固定检测器，是交通流预测最常见的数据结构。\textbf{订单}单位中每行是一笔出行订单，包含起点、终点、出发时间、行程时长和费用，来自网约车或公共交通系统，适合分析出行需求和路径选择。\textbf{事故}单位中每行是一起事故记录，包含发生时间、地点、类型和严重程度，来自警方或保险公司，适合事故频率分析和风险预测。\textbf{司机-日}单位中每行是某司机某一天的汇总指标（如接单量、行驶里程和在线时长），来自平台运营日志，适合分析驾驶行为或劳动力供给。\textbf{站点-小时}单位中每行是某公交站点某小时的乘客上下车量，来自刷卡数据的聚合，适合公交客流分布分析。

观测单位不同，所适合的特征、标签和模型假设就完全不同。以“预测拥堵”为例：如果观测单位是路段-5分钟，那么特征应包含历史流量、时段编码和天气；如果观测单位是订单，那么特征应包含出发地、目的地和出发时间；两个问题形式上都叫“拥堵预测”，但建模框架完全不同。

\subsection{聚合粒度的选择}

许多数据集来自原始高频记录，需要经过时间或空间聚合才能用于建模。聚合粒度的选择是一个重要的建模决策，而非纯粹的技术细节。

以时间粒度为例，考虑将5分钟级检测器数据聚合到不同粒度的影响：5分钟粒度保留最多的时序细节，适合短期预测，但数据量大，噪声较高，缺失值的影响也更显著；1小时粒度平滑了短期波动，适合分析日内规律，但掩盖了高峰期内部的动态变化（例如早高峰中的峰值时刻）；1天粒度只保留日均水平，适合分析工作日/周末规律或季节趋势，但失去了全部日内信息。

粒度越粗，信息损失越大，但数据也越平稳、缺失率越低、计算成本越小。正确的粒度选择取决于你的预测目标所要求的时间分辨率：如果下游应用是交通信号控制，5分钟粒度是合理的；如果是城市规划评估，日均数据可能已经足够。

空间聚合同样需要谨慎。将若干检测器的数据聚合到“某路段”或“某区域”时，必须考虑所选范围是否具有同质性：如果路段内部存在明显的流量分层（例如靠近出口匝道的位置流量显著更高），强行合并会掩盖重要的空间异质性。

\subsection{时空聚合的陷阱}

聚合操作存在两个经典陷阱，在交通数据分析中尤为常见。

\textbf{辛普森悖论}（Simpson's Paradox）：在分组数据中成立的规律，在合并数据后可能消失甚至反转。一个典型例子：分别分析工作日早高峰和非高峰时段，两个时段的平均速度都与降雨呈负相关（下雨时速度更低）；但如果不控制时段直接混合，由于降雨更多发生在非高峰时段（统计上的巧合），混合后的数据中可能出现“降雨与速度无显著关系”的错误结论。

\textbf{生态学谬误}（Ecological Fallacy）：以聚合单位得出的结论，不能直接推论到个体层面。例如，某城市交通区域的平均行程时间与该区域的人口密度呈正相关，并不意味着在该区域内，人口越密集的路段行程时间一定越长——区域层面的相关性不能推断路段层面的因果关系。

这两个陷阱的共同根源是：聚合操作是有损的，损失的信息有时恰恰是关键变量。在进行任何跨层次推断时，需要明确说明当前分析的观测单位，并意识到结论的适用范围。

\subsection{在Prompt中描述观测单位}

向 AI 描述数据时，明确观测单位是最重要的信息之一。以下是一个清晰的观测单位描述示例：

\begin{quote}
数据集中每一行代表某个检测器在某个5分钟时间窗口内的汇总记录。主键为（detector\_id, timestamp）的组合，其中 timestamp 为该时间窗口的起始时间。数据的时间范围为2023年1月1日至12月31日，共20个检测器，每天每个检测器有288条记录（24小时 × 12个5分钟窗口）。
\end{quote}

对比之下，只告诉 AI“有20个检测器的全年数据”而不说明每行代表什么，AI 可能错误地将每行理解为一个检测器的全天汇总，或者一辆车经过检测器的单次记录，导致后续分析的方向完全偏离。

% ----------------------------------------------------------------
\section{标签是如何产生的}
\label{sec:label}
% ----------------------------------------------------------------

在监督学习任务中，标签（label）是模型的学习目标。许多使用者习惯性地把标签当作“正确答案”，但这种理解掩盖了一个关键事实：\textbf{标签不是客观事实，而是人为定义的操作化结果}。标签的质量直接决定了模型能学到什么。

\subsection{标签的来源类型}

交通工程中的标签通常来自几类不同的来源，每类来源都有其特有的质量问题。

\textbf{传感器阈值判断}根据传感器读数与预设阈值的比较，自动生成二分类标签，例如“当路段平均速度低于20 km/h时，标记为拥堵”。这类方式的问题在于阈值的选择是主观的，在不同道路等级、不同城市的含义完全不同；传感器本身的误差会直接传递到标签；阈值附近的记录（例如速度为21 km/h的路段）被标记为非拥堵，但实际行驶感受可能与速度为5 km/h的拥堵差异并不显著。

\textbf{警方或机构记录}提供事故严重程度、违规类型等标签。这类标签的问题在于记录质量受执法能力和记录规范的影响，漏报率在轻微事故中可能高达60\%以上；同一机构在不同年份的记录标准可能发生变化，例如“重伤”的医学定义修订后，历史数据和新数据的可比性会下降；地理信息通常为文字描述，匹配到路网时误差不可避免。

\textbf{人工标注}为视频、图像或文本数据提供分类标签。这类方式的问题在于标注员间一致性（inter-annotator agreement）往往低于预期，尤其是对于边界模糊的类别；标注指南的清晰度直接影响标注质量，而指南通常在项目进行中才逐步完善，导致早期标注与晚期标注使用了不同的隐含规则；标注员的工作状态（疲劳、动机）会影响标注质量，但这种影响几乎无法事后量化。

\textbf{规则派生}从原始数据通过业务规则推导出标签，例如将公交车到站时间与计划表对比判断“准点”或“晚点”，或根据车载传感器的加速度序列判断驾驶行为类型。这类方式的问题在于规则本身可能不完整或存在边界情况，而这些边界情况在大规模数据中会被大量触发；规则的制定者通常针对“正常情况”设计规则，对于极端场景（暴雪、重大事故）的处理可能是未定义的。

\textbf{平台机制派生}标签（如网约车订单的“爽约”、“投诉”）来自平台系统，其生成逻辑受到平台规则的影响。一个典型例子：平台改变了“司机爽约”的判定时间窗口（从等待10分钟改为5分钟），导致同一司机行为在系统调整前后被赋予不同标签，而历史数据和新数据在标签含义上已不再一致。

\subsection{标签定义对模型的影响}

不同的标签定义，训练出的模型回答的是不同的问题。“拥堵”可以被定义为：

\begin{enumerate}
  \item 路段平均速度低于20 km/h；
  \item 行程时间比自由流时间延误超过50\%；
  \item 交通量超过通行能力的80\%；
  \item 驾驶员自评感受（来自调查问卷）。
\end{enumerate}

这四种定义会产生四份不同的标签序列：同一条路段、同一个时刻，在不同定义下可能分别被标记为拥堵和非拥堵。基于这四份标签分别训练的模型，其输出在业务上代表完全不同的含义。如果一个工程师基于定义一训练了模型，另一个工程师基于定义二在同一路段评估模型性能，两人得到的指标不具有可比性，但这种不可比性在标签不透明时很难被发现。

\subsection{向AI描述标签时应说明的信息}

在向 AI 提供含有标签的数据集时，至少应说明以下内容：标签的字段名和取值（例如字段名为 \texttt{congestion}，取值为0或1，1代表拥堵）；标签的操作化定义（例如当5分钟平均速度低于20 km/h时标记为1，否则为0）；标签的生成方式（自动阈值判断、人工标注、警方记录等）；标签的已知质量问题（例如轻微事故漏报率估计在40\%—60\%之间、某年份之前的严重程度分级标准与当前不同）；以及标签的时间范围内是否发生过定义变更。

不向 AI 说明标签含义，它会将标签字段视为需要预测的数值，而不理解该数值在业务上代表什么。这会导致 AI 选择不适合的损失函数（例如对二分类标签使用均方误差）、使用不适合的评价指标，或者无法正确处理标签的不平衡问题。

% ----------------------------------------------------------------
\section{数据的质量问题：人必须亲自判断的部分}
% ----------------------------------------------------------------

数据质量问题是每个实际数据项目都无法绕过的现实。AI 可以帮助检测质量问题的表现（例如统计缺失率、绘制异常值分布），但它无法判断这些问题在你的具体场景下意味着什么，以及应当如何处理。这种判断需要领域知识，而领域知识只有你才有。

\subsection{缺失值的含义不唯一}

缺失值（missing value）的危险在于：相同的表现（字段为空或为 NaN）在不同场景下代表完全不同的含义，而处理方法也因此截然不同。

以交通检测器数据为例，同一字段的缺失可能来自截然不同的原因。\textbf{传感器故障}是指硬件损坏或通信中断导致数据未传输，这类缺失相对于交通状态而言是随机的，可以考虑用插值或前向填充处理，但需注意长时间缺失段不应盲目插值。与之不同，\textbf{真实零流量}是指深夜低流量时段，路段确实没有车辆通过，检测器返回零而非缺失，如果把这类零值当作缺失处理，会系统性地删除低流量时段的有效数据，破坏日内流量的自然规律。\textbf{通信延迟后补传}是指数据实时传输失败但后来通过补传机制恢复，在历史数据中这类缺失已被填补，但在实时预测场景中该时刻数据确实不可用，预测时不能使用事后补传的数据作为特征。\textbf{设计性缺失}则是指某类数据本来就不采集某些时段或路段，例如某些检测器在节假日会被手动关闭以节省存储，这类缺失是结构性的，插值没有意义。

上述四种缺失，AI 在没有额外说明的情况下无法区分。它通常会套用最常见的默认策略（例如均值填充或线性插值），这在某些情况下是合适的，在另一些情况下会引入严重偏差。你必须明确告诉 AI 你数据中缺失值的机制，以及对应的处理规则。

\subsection{异常值的两面性}

异常值（outlier）的处理比缺失值更为复杂，因为它面临一个根本性的不确定性：\textbf{这个极端值是测量错误，还是真实发生的极端事件？}

在交通数据中，这两种情况都存在。检测器漂移可以导致速度读数持续偏高（例如某检测器长期报告100 km/h以上的速度，而该路段限速60 km/h），这是测量错误，应当识别并标记为无效。另一方面，重大交通事故可能导致某路段流量在几分钟内骤降至近零随后恢复，这是真实的极端事件，在事故研究中恰恰是最有价值的数据点，如果被当作异常值删除，会严重损害事故影响分析。此外，特殊活动（演唱会散场、体育赛事结束）导致周边路段短时间内出现极高流量，对于大多数常规预测任务来说是“离群点”，但在特殊事件分析中却是核心数据。

区分这两类情况需要领域知识：了解路段的实际限速和正常流量范围、知道数据采集时间段内是否有已知的特殊事件、了解该检测器的历史维护记录。这些信息存在于领域专家的头脑中，而不在数据集里。

一个可操作的原则是：\textbf{在删除或修改任何异常值之前，先问自己“这个值在现实中能否发生”，再问“如果真实发生，我的分析目标是否需要保留它”}。

\subsection{时间对齐问题}

当一个分析需要整合来自不同来源的数据时，时间对齐（temporal alignment）是一个必须明确处理的问题，而不是可以留给 AI 自动推断的细节。

常见的时间对齐陷阱有以下几类。\textbf{时区不一致}方面，部分检测器数据以 UTC 存储，部分以本地时间存储，合并时未做转换，导致早高峰时段的记录在合并数据中对应凌晨时段。\textbf{采样频率不匹配}方面，将5分钟粒度的检测器数据与小时粒度的气象数据合并时，需要明确决策：是将气象数据复制到每个5分钟时间步，还是将检测器数据聚合到小时，不同选择对后续分析的含义不同。\textbf{时间戳含义不统一}方面，某些系统记录的是时间窗口的起始时间（例如7:00代表7:00—7:05），另一些系统记录的是结束时间（7:05代表7:00—7:05），还有一些系统记录的是采集时刻（7:03代表该时刻的瞬时值），将这三类时间戳对齐时，如果不加说明，容易引入半个到一个时间窗口的系统性错位。\textbf{夏令时处理}方面，在有夏令时的地区，每年会有一个25小时的天和一个23小时的天，如果时间戳未被正确转换，这两天的数据会出现一小时的重叠或缺失，导致时序特征（如“小时编码”）在这两天产生错误。

\subsection{空间对齐问题}

与时间对齐类似，空间对齐（spatial alignment）问题在多源数据融合时同样常见且危险。\textbf{路段编号不一致}是指交通管理部门使用自己的路段编码体系，地图供应商使用另一套，道路维护部门又用第三套，不同数据集的路段编号可能指向不同的空间范围，强行用编号合并会导致严重的空间错配。\textbf{坐标系不一致}是指国内地图数据常见 GCJ-02（火星坐标）、BD-09（百度坐标）和 WGS-84（GPS 国际标准）三种坐标系，未经转换直接使用会导致几十米到几百米的系统性偏移，在城市道路精度要求较高的分析中会产生严重错误。\textbf{地图版本不一致}是指同一条路在不同年份的地图数据中，路段划分方式可能发生变化（例如某路段被一分为二，或两条支路被合并），将不同年份的数据合并时，历史数据和当前数据的路段对应关系可能已经不成立。

\subsection{数据偏差}

除了上述数据质量问题，数据偏差（data bias）是一种更深层、更难察觉的问题。它不表现为明显的缺失或异常，而是数据的采集设计本身就只能观测到现实的一个子集，且这个子集不是随机的。

\textbf{位置偏差}是指固定检测器的安装位置通常由历史交通管理需求决定，而不是随机布设，主干道、交叉口和高速公路的检测器密度远高于支路和居民区道路。用这些数据训练的模型，在预测支路交通状态时性能会系统性地更差，但这种差异在测试集上可能看不出来（如果测试集也来自同样的检测器）。\textbf{人群偏差}是指手机导航 APP 的数据只来自使用该 APP 的用户群体，这一群体在年龄、收入、出行目的等方面与全体出行者存在系统性差异，基于 APP 数据得出的出行规律，不能不加修正地推广到所有出行群体。\textbf{时间偏差}是指网约车数据只记录了有叫车需求且被满足的出行，未被满足的需求（用户叫车后放弃等待）和不选择网约车的出行都不在数据集中，如果用这份数据估计出行需求，会系统性地低估高峰时段的潜在需求（因为高峰时段供不应求，部分需求未被满足）。

数据偏差的根本问题在于：模型学到的是“数据中呈现的规律”，而不是“现实世界的规律”。当数据的采集过程系统性地遗漏了某些情况时，模型在这些情况下的表现不是随机的不准确，而是有方向性的偏差。

\subsection{为什么这些问题AI无法自动处理}

以上每一类问题，都需要在数据之外的知识来做判断：知道检测器的安装位置和维护历史，知道数据采集系统在哪个时间段升级过，知道平台在哪个月份修改过派单规则，知道某个路段在哪年进行了改扩建。

AI 看到的只是一张表格——无论这些背景知识多么重要，如果不被明确地写入 Prompt，AI 都无从知晓。因此，\textbf{数据质量的判断必须是人的工作，AI 最多只能协助生成初步的质量检查代码，以及对你提供的质量信息做出响应}。

% ----------------------------------------------------------------
\section{向AI描述数据：构建数据上下文}
% ----------------------------------------------------------------

第二章已经讨论了“上下文”作为 Prompt 五要素之一的重要性。针对数据分析任务，上下文的核心是数据上下文（data context）——向 AI 解释你的数据是什么，使 AI 能够在正确的前提下生成代码和分析。

\subsection{数据上下文的六个要素}

一份完整的数据上下文应当覆盖以下六个方面。\textbf{来源}说明数据从哪里来，通过什么采集方式获得，例如“数据来自某市交通管理局部署的感应线圈检测器，原始数据为每分钟一条记录，本数据集已汇总为5分钟粒度”。\textbf{观测单位}说明每行代表什么，例如“每行代表某个检测器在某个5分钟时间窗口内的汇总，主键为 (detector\_id, window\_start) 的组合”。\textbf{字段语义}说明每个字段的含义、单位和取值范围，特别是与通用含义不同之处，例如“speed 字段为该时间窗口内所有过车的平均速度，单位为 km/h；speed = 0 代表检测器离线，不是真实零速度；flow 字段为该时间窗口内的总过车数，单位为辆；flow = $-1$ 代表检测器故障，应视为缺失值”。\textbf{时空范围}说明数据覆盖的时间范围、空间范围和采样规律，例如“时间范围为2023年1月1日至12月31日，共20个检测器，分布在城市东部快速路沿线；每天每个检测器理论上应有288条记录，但实际因故障存在缺失”。\textbf{质量情况}说明已知的缺失率、异常值比例和特殊情况，例如“全数据集整体缺失率约8\%，其中3号检测器在6月至8月间因施工停用，该时段缺失率为100\%；speed 字段有约2\%的记录超过限速（60 km/h），其中部分来自轻型摩托车，部分疑为传感器错误”。\textbf{标签定义}说明如果有标签字段，该字段的含义、生成方式和已知质量问题（参见第\ref{sec:label}节）。

\subsection{好的数据上下文与坏的数据上下文}

以一个典型任务为例：请 AI 生成代码，分析检测器数据中的缺失模式。

\textbf{坏的数据上下文：}

\begin{quote}
我有一份交通检测器数据，请分析缺失值情况。
\end{quote}

这条 Prompt 几乎不包含任何有用的信息。AI 不知道哪些字段需要分析、缺失在数据中如何表示（NaN、-1还是0）、哪些字段的缺失是需要重点关注的。它会生成一份通用的缺失值分析代码，统计每列的 NaN 比例——这个结果对于 speed = 0 代表缺失的情况是完全错误的。

\textbf{好的数据上下文：}

\begin{quote}
我有一份交通检测器数据，文件名为 \texttt{detector\_2023.csv}，字段包括：
\begin{itemize}
  \item \texttt{timestamp}：时间窗口起始时间，格式为 \texttt{YYYY-MM-DD HH:MM:SS}，时区为北京时间（UTC+8）；
  \item \texttt{detector\_id}：检测器编号，整数，范围1—20；
  \item \texttt{speed}：平均速度，单位 km/h；\textbf{speed = 0 代表检测器离线，不是真实零速度，应视为缺失值处理}；
  \item \texttt{flow}：5分钟过车数，单位辆；\textbf{flow = $-1$ 代表检测器故障，应视为缺失值}；
  \item \texttt{occupancy}：占有率，取值0—1；该字段无特殊约定，NaN即为缺失。
\end{itemize}
已知3号检测器在2023年6月1日至8月31日因道路施工停用，该时段的数据全部缺失。

请分析以下缺失情况：（1）以正确方式处理各字段的缺失值编码（speed = 0 和 flow = $-1$ 应先转为 NaN）；（2）统计每个检测器每个字段的有效数据率（非缺失比例）；（3）按月份统计整体缺失率，并标注3号检测器施工期间的数据情况；（4）输出 Markdown 格式的摘要表，并绘制检测器编号×月份的缺失率热力图，颜色越深代表缺失率越高。
\end{quote}

这条 Prompt 向 AI 提供了足够的信息：哪些值代表缺失（speed = 0，flow = $-1$）、每个字段的语义、已知的特殊情况（3号检测器停用）以及具体分析任务。AI 能够在此基础上生成真正有用的代码，而不必靠猜测填补信息空白。

\subsection{数据字典的作用}

在规模较大的项目中，数据字典（data dictionary）是维护字段语义的标准工具。数据字典通常是一份表格，记录每个字段的名称、类型、单位、取值范围、含义和特殊说明。

数据字典对 AI 协作的价值在于：它可以直接被复制到 Prompt 的上下文部分，作为字段语义的权威说明。这比每次重新手写字段说明更高效，也更不容易遗漏关键信息。

一个简洁的数据字典示例如下：

\begin{quote}
\textbf{detector\_2023.csv 数据字典}

\begin{itemize}
  \item \texttt{timestamp}（字符串）：时间窗口起始时间，格式 YYYY-MM-DD HH:MM:SS，UTC+8。
  \item \texttt{detector\_id}（整数）：检测器编号，1—20，对应城市东部快速路沿线20个感应线圈。
  \item \texttt{speed}（浮点数，km/h）：5分钟内过车平均速度。\textbf{特殊值：0 = 检测器离线（非真实零速）}，NaN = 数据未传输。
  \item \texttt{flow}（整数，辆/5分钟）：5分钟内过车总数。\textbf{特殊值：$-1$ = 检测器故障}。
  \item \texttt{occupancy}（浮点数，0—1）：车辆占用检测区的时间比例，NaN = 传感器异常。
\end{itemize}
\end{quote}

\subsection{常见错误：只提供字段名，不提供字段含义}

最常见的数据上下文错误，是在 Prompt 中只告诉 AI 字段名列表，而不说明每个字段的含义。

\begin{quote}
我的数据有以下字段：timestamp, detector\_id, speed, flow, occupancy, lane\_count, heavy\_vehicle\_ratio。请分析这份数据。
\end{quote}

对于 \texttt{speed} 和 \texttt{flow} 这类字段名，AI 可以合理猜测含义。但对于 \texttt{heavy\_vehicle\_ratio}，AI 无法知道它是百分比（0—100）还是比例（0—1），也不知道“重型车辆”的分类标准（是按轴数、车长还是车重？），更不知道该字段的缺失值编码方式。凭猜测生成的分析代码，在这些细节上可能完全错误，而错误在输出结果中往往不会以明显方式暴露。

% ----------------------------------------------------------------
\section{与AI协作进行探索性数据分析}

探索性数据分析并不是在正式建模前机械生成一组图表，而是通过可视化、变换和反复提问发现数据结构、异常及假设。Tukey 将探索与验证区分开来，强调先让数据暴露值得进一步检验的问题\cite{tukey1977exploratory}。AI 可以加速生成代码和提出观察角度，但使用者仍需判断某个模式究竟来自真实现象、采集制度，还是处理错误。
\label{sec:eda}
% ----------------------------------------------------------------

探索性数据分析（Exploratory Data Analysis，EDA）是数据科学工作流中的核心环节。它的目的不是产出最终结论，而是在建模之前充分理解数据的结构、质量和规律，为后续决策提供依据。AI 在 EDA 的代码生成环节有显著的效率优势，但 EDA 的核心价值——发现有意义的模式并结合领域知识做出判断——仍然需要人来完成。

\subsection{EDA的目的}

一次完整的 EDA 至少应当回答以下几类问题：

\begin{enumerate}
  \item \textbf{数据结构是否符合预期？}行数、列数、时间范围、检测器数量——这些基本指标与数据说明文档是否一致？
  \item \textbf{数据质量如何？}各字段的缺失率、异常值比例、特殊值的分布情况是否在可接受范围内？
  \item \textbf{数据分布是否合理？}各数值字段的均值、中位数、分布形状和极值是否符合领域常识？
  \item \textbf{存在哪些时序和空间规律？}交通数据通常具有强烈的日内规律（早晚高峰）和周内规律（工作日与周末），是否能在数据中观察到这些规律？
  \item \textbf{变量之间的关系如何？}与标签相关的特征是否表现出预期的关联性？不同变量之间是否存在多重共线性？
\end{enumerate}

EDA 的发现直接影响后续的数据清洗策略（第五章）和特征工程决策（第六章）。一次不充分的 EDA，意味着在不完整的信息基础上做出数据处理决策，后患无穷。

\subsection{AI适合做什么}

在 EDA 阶段，AI 能够高效完成的工作集中在任务明确、有标准实现方式且输出可被直接验证的环节：生成基本统计摘要（行数、列数、各字段的均值、标准差、最大值、最小值、分位数）；统计缺失率、计算各字段的非空记录数；绘制数值字段的分布直方图和箱线图；计算变量间的相关性矩阵并可视化为热力图；检测缺失值的空间模式（哪些检测器缺失率更高）和时间模式（哪些时段缺失更集中）；生成时序图，按检测器或按时间粒度展示数据变化；绘制工作日与周末的对比分布图。AI 生成的代码在大多数情况下可以直接运行，即使需要调整也有明确的方向。

\subsection{人必须做什么}

AI 无法替代的 EDA 工作体现在几个关键方面。\textbf{解读异常现象}方面，图表中出现的异常模式需要结合领域知识解释——例如，某检测器在春节期间流量骤降，AI 可以检测到这个模式，但只有了解中国节假日出行规律的人才能判断这是正常现象，而不是数据错误。\textbf{判断缺失机制}方面，AI 可以告诉你缺失率是多少，但判断缺失是“传感器随机故障”还是“某时段数据从未被采集”，需要你了解数据采集系统的工作方式。\textbf{验证AI的结论}方面，AI 有时会给出看起来合理、实则错误的发现——例如，AI 可能“发现”两个变量之间存在强相关性，但这个相关性实际上来自时间趋势的共同驱动，而不是两个变量之间的真实关联（这是一种常见的虚假相关），识别这类问题需要统计知识和领域直觉。\textbf{形成建模假设}方面，EDA 的最终目的是为建模提供依据，哪些特征可能有预测力、数据的时序结构是否要求特殊的建模方法，这些判断需要综合 EDA 发现和领域知识，不能由 AI 代为做出。

\subsection{EDA的推荐流程}

一次完整 EDA 可以按照以下顺序推进，每个步骤对应一次或多次 AI 协作：

\begin{enumerate}
  \item \textbf{数据概览：}读取数据，输出基本统计信息（行数、列数、时间范围、各字段非空数量）。\textit{检查：实际行数与预期是否一致；字段类型是否被正确识别（时间戳是否被解析为日期时间格式）。}
  \item \textbf{字段分布：}对每个数值字段绘制分布图（直方图、箱线图），输出均值、中位数、标准差和分位数。\textit{检查：数值范围是否合理；是否存在明显的双峰分布或长尾，且这种形状是否能被领域解释。}
  \item \textbf{缺失与异常：}统计每个字段的缺失率（注意特殊值编码），输出缺失值的时间分布和检测器分布。检测数值异常（远超合理范围的记录）。\textit{检查：缺失模式是否与已知的传感器故障时段吻合；异常值是否集中在特定检测器或特定时段。}
  \item \textbf{时序与空间模式：}绘制典型工作日和周末的日内变化曲线，检验早晚高峰是否存在；按检测器分组绘制流量和速度的分布对比图。\textit{检查：日内规律是否符合交通常识（早高峰7—9时，晚高峰17—19时）；不同检测器之间的规律是否存在合理的空间差异。}
  \item \textbf{变量关系：}计算各特征与标签之间的相关性，绘制相关性矩阵热力图；对强相关的变量对绘制散点图，确认关系的形态。\textit{检查：相关性的方向是否符合领域预期；强相关是否可能来自第三变量的混淆。}
\end{enumerate}

\subsection{用Prompt引导AI做EDA：一个完整示例}

以下是一条完整的 EDA 任务 Prompt，展示如何提供充分的数据上下文并引导 AI 生成有用的分析代码：

\begin{quote}
\textbf{数据说明：}

我有一份交通检测器数据，文件名为 \texttt{detector\_2023.csv}，包含以下字段：
\begin{itemize}
  \item \texttt{timestamp}（字符串）：时间窗口起始时间，格式 \texttt{YYYY-MM-DD HH:MM:SS}，北京时间（UTC+8）；
  \item \texttt{detector\_id}（整数，1—20）：检测器编号；
  \item \texttt{speed}（浮点数，km/h）：5分钟内过车平均速度；speed = 0 代表检测器离线（非真实零速），应转为 NaN；
  \item \texttt{flow}（整数，辆/5分钟）：5分钟内过车总数；flow = $-1$ 代表检测器故障，应转为 NaN；
  \item \texttt{occupancy}（浮点数，0—1）：车辆占用检测区的时间比例，NaN 即缺失。
\end{itemize}

已知情况：3号检测器在2023年6月1日至8月31日因道路施工停用。

\textbf{任务：}

请完成以下探索性数据分析，生成一段可直接运行的 Python 脚本（Python 3.10，使用 pandas、matplotlib 和 seaborn，不引入其他第三方库）：

\begin{enumerate}
  \item 读取数据，将 \texttt{timestamp} 解析为 datetime 类型，将 speed = 0 和 flow = $-1$ 转为 NaN；
  \item 输出数据基本信息：总行数、时间范围（最早和最晚时间戳）、各字段的非空记录数和有效数据率；
  \item 统计每个检测器每个字段的缺失率，并以检测器编号为纵轴、字段为横轴绘制缺失率热力图，颜色越深代表缺失率越高；在图注中特别标注3号检测器；
  \item 绘制 speed 和 flow 字段的分布直方图（各一张），在图上标注均值和中位数；
  \item 按星期几（周一至周日）分组，绘制工作日（周一至周五）与周末（周六至周日）的小时平均速度曲线对比图，横轴为小时（0—23），纵轴为平均速度；
  \item 计算 speed、flow 和 occupancy 三个字段之间的相关性矩阵（仅使用三个字段均非空的记录），并绘制带数值标注的热力图。
\end{enumerate}

\textbf{输出要求：}

每个分析步骤前添加一行注释说明该步骤的目的；图表保存为 PNG 文件，文件名与步骤编号对应；在所有分析完成后，输出一段 Markdown 格式的摘要，包含数据概况和各步骤发现的主要模式，其中“发现的主要模式”部分留空（标记为“[待人工填写]”），以提醒使用者需要结合领域知识自行解读。
\end{quote}

这条 Prompt 的最后一项要求——将发现的主要模式留空等待人工填写——体现了一个重要原则：AI 生成的统计图表需要人的解读，而解读需要领域知识，不应由 AI 代劳。

\subsection{警惕AI的自动结论}

当 AI 在 EDA 代码之外主动给出“发现”和“结论”时，需要保持批判性审视。常见的风险体现在几个方面：\textbf{混淆相关与因果}，AI 可能报告“降雨量与路段速度呈负相关，说明降雨导致速度下降”，但“导致”是因果推断，不能仅从相关性中得出；\textbf{忽视混淆变量}，AI 看到的相关性可能来自第三个变量的共同影响，而不是两个变量之间的直接关系；\textbf{过度解读数值}，AI 可能将统计上显著但业务上微不足道的差异报告为“重要发现”，而判断一个差异是否具有实际意义需要领域知识；\textbf{基于错误假设的推断}，如果 AI 对数据的某个字段的含义理解有误（例如把检测器离线的 speed = 0 当作真实低速），其后续所有分析都会在错误的基础上进行，且错误不会在代码执行时报错。

原则是：\textbf{让 AI 负责生成图表和统计数值，由人来做模式识别和意义解读}。AI 生成的结论可以作为假设的来源，但不应不经验证地作为分析结论采纳。

% ----------------------------------------------------------------
\section{实战练习}
% ----------------------------------------------------------------

\subsection*{练习一：识别数据类型的内生局限}

以下是三种交通数据采集方式。对每种方式，请说明：（1）它能观测到什么、不能观测到什么；（2）它的主要内生局限是什么；（3）如果用这种数据研究“城市拥堵的时空分布规律”，可能产生什么系统性偏差。

\begin{enumerate}
  \item 城市主干道的固定感应线圈检测器网络（覆盖率约40\%的路段）；
  \item 某网约车平台的行程 GPS 轨迹数据（平台市场占有率约30\%）；
  \item 手机信令数据（由移动运营商提供，覆盖该城市约70\%的手机用户）。
\end{enumerate}

\subsection*{练习二：分析观测单位与建模设计的关系}

某研究团队希望预测城市公交系统的乘客到站需求，以便提前调度车辆。他们拥有以下两份数据：

\begin{itemize}
  \item 数据集 A：每行代表一笔刷卡记录，字段包括卡号、上车站点、上车时间、下车站点（部分缺失）；
  \item 数据集 B：每行代表某站点某小时的乘客上车量，字段包括站点编号、日期、小时、上车人数。
\end{itemize}

请回答：（1）这两份数据集的观测单位分别是什么？（2）如果直接用数据集 A 建模，需要先进行什么转换？（3）数据集 B 在空间粒度上可能存在什么信息损失？（4）请写一条向 AI 描述数据集 B 的 Prompt 片段（只需数据上下文部分，不需要写出完整任务）。

\subsection*{练习三：诊断标签质量问题}

某研究项目使用交通事故记录数据，目标是训练一个预测路段事故风险的分类模型，标签为“高风险”（过去三年该路段事故数量 $\geq$ 3 次）或“低风险”（$<$ 3 次）。

请分析以下可能影响标签质量的问题：

\begin{enumerate}
  \item 事故记录来源于不同派出所，各所的记录详细程度不同，部分轻微事故未被记录——这会如何影响“高风险”标签的分布？
  \item 三年观察窗口期内，某路段进行了大规模改造，事故率在改造后显著下降——仍然使用三年总量作为标签，会引入什么问题？
  \item 该城市在数据时段中期安装了新的交通安全设施，导致有设施区域的事故率下降——这会如何影响标签的空间分布，以及用该数据训练的模型的泛化能力？
\end{enumerate}

如果你要向 AI 描述这份标签，应当包含哪些关键信息？请写出一段数据上下文说明。

\subsection*{练习四：设计一个完整的EDA Prompt}

选择一份你实际使用过或熟悉的数据集（可以是公开数据集，例如 PeMS 加州道路检测器数据、GTFS 格式的公交数据，或 NYC TLC 出租车行程数据）。

请完成以下工作：

\begin{enumerate}
  \item 用六要素框架（来源、观测单位、字段语义、时空范围、质量情况、标签定义）写出这份数据集的完整上下文说明，字数不少于200字；
  \item 基于你对这份数据的了解，列出这份数据中你认为最重要的三个质量问题，说明每个问题的表现形式、可能原因和处理方向；
  \item 按照本章介绍的 EDA 流程，设计一条完整的 EDA Prompt，要求：包含完整的数据上下文，明确说明哪些特殊值需要被识别，列出至少5个具体的分析任务，并在输出要求中注明 AI 不应代替人工填写“发现的意义和解读”。
\end{enumerate}





